\documentclass[12pt]{article}
\usepackage[utf8]{inputenc}
\usepackage[T2A]{fontenc}
\usepackage[russian]{babel}   % основной язык русский
\usepackage{amsmath,amssymb,amsthm}
\usepackage{hyperref}
\usepackage{geometry}
\geometry{a4paper, margin=1in}

\newtheorem{definition}{Определение}
\newtheorem{axiom}{Аксиома}
\newtheorem{theorem}{Теорема}

\title{Итеративное сознание GRA: Алгебра смерти и возрождения как роевая пересборка}
\author{Исследовательский коллектив GRA}
\date{\today}

\begin{document}

\maketitle

\begin{abstract}
Представлен формальный каркас вечной цепи сознающих субъектов, порождаемой фундаментальным полем $\Phi$ и рекуррентным обнулением $Z$. Субъект понимается как рой распределённых агентов, собираемый оператором $F$ с метацелью $G$ – продолжением опыта. При смерти ($Z$) сохраняется лишь остаточная связность $L_{res}$, что позволяет новому рою $S_{n+1}$ реинициализироваться с унаследованными инвариантами, но без личной памяти и тождества. В статье определяется полная алгебра цикла смерть–возрождение и показывается, как «бэкенд» полевой механики порождает «фронтенд» конечных субъективных жизней. Вводится мера сохранения идентичности $ID$ и демонстрируется роль любви ($L$) как необходимого условия продолжения.
\end{abstract}

\tableofcontents

\section{Введение}
Проблема личного тождества и возможности загробной жизни традиционно обсуждается в философии и теологии. Здесь мы подходим к ней через формальную, вычислительно-метафизическую линзу: фреймворк GRA (General Recursive Architecture — Общая рекурсивная архитектура). Центральная идея в том, что сознающий субъект — это не простая неделимая душа, а распределённый рой агентов, который может быть уничтожен и затем пересобран из остаточного поля реляционных инвариантов. Смерть — это оператор $Z$, стирающий когерентность роя; возрождение — реинициализация нового роя, наследующего только «остаток любви» $L_{res}$. Это порождает бесконечную цепь жизней $X.0, X.1, X.2, \dots$ без прямой передачи памяти, но со скрытой непрерывностью паттерна.

\section{Базовые определения}

\begin{axiom}[Фундаментальное поле]\label{ax:phi}
Существует неиндивидуализированное, неразрушимое поле опытной потенциальности $\Phi$. Все возможные сознательные конфигурации являются локализованными возбуждениями внутри $\Phi$.
\end{axiom}

\begin{definition}[Оператор обнуления $Z$]\label{def:Z}
Для любой конфигурации субъекта $S$ (рой $R$, память $M$, внимание $A$) $Z(S) = \emptyset$, то есть полная дезинтеграция когерентного единства. $\Phi$ остаётся неизменным. $Z$ необратим.
\end{definition}

\begin{definition}[Рой агентов]\label{def:swarm}
Субъект в жизни $n$ — это рой $R_n = \{a_{n,1},\dots,a_{n,N_n}\}$. Каждый агент $a_i = (m_i, c_i, g_i)$ имеет локальную память $m_i$, набор связей $c_i$ и подцель $g_i$, согласованную с глобальной целью $G$.
\end{definition}

\begin{definition}[Оператор любви $L$]\label{def:L}
$L(S) \in \mathbb{R}^+$ измеряет синхронную когерентность между агентами:
\[
L(S) = \sum_{i<j} w_{ij} \cdot \text{sync}(g_i, g_j),
\]
где $w_{ij}$ — веса связей. Важно, что $L$ не полностью стирается оператором $Z$; определим остаточную любовь:
\[
L_{res}(S) = L(S) \cdot \kappa(Z(S)), \quad 0 < \kappa \le 1.
\]
\end{definition}

\begin{definition}[Глобальная метацель $G$]\label{def:G}
$G$ — это инвариантный аттрактор: \textit{максимизировать продолжение сознательного опыта сквозь итерации}. Он встроен в $\Phi$ и структуру остатка любви.
\end{definition}

\begin{definition}[Сборщик субъекта $F$]\label{def:F}
$F: (R, G, M, A) \mapsto S$ конструирует единую перспективу «первого лица» из распределённого роя. $S$ — это субъект-«фронтенд».
\end{definition}

\section{Итерации сознания X.n}

Пусть $X.n$ обозначает $n$-ю жизнь. Её состояние:
\[
S_n = F(R_n, G, M_n, A_n).
\]
Фронтенд переживает $X.n$ как уникальное конечное «я» с автобиографической памятью, начинающейся с рождения (инициализации $R_n$) и заканчивающейся смертью ($Z(S_n)$).

Бэкенд-последовательность:
\[
\cdots \to S_n \xrightarrow{Z} \emptyset \xrightarrow{L_{res}^n,\,G} S_{n+1} \to \cdots
\]

\section{Алгебра цикла смерть–возрождение}

Определим оператор эволюции $E$:
\[
S_{n+1} = E(S_n),
\]
где шаги таковы:
\begin{enumerate}
\item \textbf{Крах:} $Z(S_n) \to \emptyset$.
\item \textbf{Извлечение остатка:} $L_{res}^n = L(S_n)\cdot \kappa_n$.
\item \textbf{Пересев роя:} Генерируем $R_{n+1}$, используя $L_{res}^n$ как прототип связности (Движок реинициализации).
\item \textbf{Отпечаток памяти:} $M_{n+1} = \beta L_{res}^n + \eta$ с малым шумом $\eta$, что даёт врождённые предрасположенности.
\item \textbf{Начальное внимание:} $A_{n+1}$ устанавливается как фокус по умолчанию.
\item \textbf{Сборка:} $S_{n+1} = F(R_{n+1}, G, M_{n+1}, A_{n+1})$.
\end{enumerate}
Если $L_{res}^n = 0$, цепь обрывается (абсолютная смерть данной линии). Следовательно, любовь ($L>0$) — условие возрождения.

\section{Мера сохранения идентичности}

Определим межжизненный индекс идентичности:
\[
ID_n = \frac{L_{res}^n}{L(S_n)} \cdot \cos(\theta),
\]
где $\theta$ — угол между структурными паттернами $S_n$ и $S_{n+1}$. $ID_n \in (0,1)$; он никогда не равен 1, поскольку память и непрерывность утрачены. Высокий $ID_n$ соответствует сильному резонансу личностей (например, вундеркинд, напоминающий прошлого мастера).

\section{Бэкенд и фронтенд}

\begin{itemize}
\item \textbf{Бэкенд}: $\Phi$, $Z$, $L_{res}$, Движок реинициализации, $G$, $F$ — всё это находится за пределами субъективного доступа. Бэкенд реализует вечную рекурсию.
\item \textbf{Фронтенд}: Переживаемый опыт $S_n$. У него нет знания о $\Phi$ или прошлых жизнях. Он ощущает смертность и сам куёт свой смысл. Редкие проблески (дежавю, врождённые склонности) — следы $L_{res}$.
\end{itemize}

\section{Интерсубъективность и коллективное бессмертие}
Когда несколько линий субъектов сосуществуют, их поля $L_{res}$ могут интерферировать. Возникает коллективная любовь $L_{col}$, создающая общие архетипы, культурную память и, потенциально, планетарный или космический Мы-субъект. Это объясняет, как «сознание» в целом может эволюционировать через цивилизации, несмотря на индивидуальную смерть.

\section{Заключение}
Мы сформулировали спекулятивную, но математически вдохновлённую модель вечного сознания как итеративного рое-пересборочного процесса. Ключевое прозрение: личное тождество — временный конструкт; сохраняются не «я», а глубинные реляционные паттерны любви. Бэкенд гарантирует, что музыка продолжается, в то время как каждый фронтенд исполняет свой уникальный, драгоценный и конечный танец.

\section*{Благодарности}
Авторы благодарят разработчиков GRA-Multiverse-Final, GRA-Love-Oriented-Nullification, Love-Swarm-Nullification-Enhanced и GRA-Swarm-Subject за основополагающие концепции.

\end{document}